\paragraph{\S105, continuation. Abelian pairings, algebraic topoi and an unfinished universal property (pp.\ 399--406).}
This continuation begins with Case~2, introduced on p.\ 398. For $B=A^o$, (49) identifies $B^\wedge_{ab}$ with the category of additive cosheaves on $A^\wedge_{ab}$ valued in $(Ab)$. The corresponding pairing (50) commutes with small direct limits in each argument, and its symmetry supplies the analogous description of $A^\wedge_{ab}$. The equivalence (51) reduces such pairings to their values on $A\times B$. Formula (52) gives the free abelian group on $\mathrm{Hom}(b,a)$ when the arguments are the abelianizations of representables; (53) extends this to $\underline Z^{(F*G)}$ for presheaves $F,G$. The insertion on p.\ 400 explicitly says that \emph{both members}, rather than only the left side, commute with the relevant direct limits (pp.\ 399--400).

Remark~1 abstracts this situation to two additive categories $\mathcal P,\mathcal Q$ with a pairing to $(Ab)$, cocontinuous in each variable, which identifies each category with the additive cosheaves on the other. The presheaf examples are abelian topoi: abelian categories with exact small filtering direct limits and a small generating subcategory. The visible revisions add the filtering qualification and clarify that this terminology does not mean a category which is simultaneously an ordinary topos and an abelian category. Replacing abelian groups by modules over a commutative ring gives the corresponding $k$-linear examples. For a noncommutative ring, the two sides use right and left modules; the proposed pairing is the abelian-group pairing followed by quotienting the resulting bimodule by the differences $sx-xs$. Grothendieck explicitly says that he has not checked that this construction gives the desired duality. He does verify the ordinary one-object example through tensor product, evaluation at the regular module, and the associated tensor--Hom equivalences. Their comparison is an additive analogue of (39) when the target category has both kinds of small limits (pp.\ 400--402).

Remark~2 singles out the existence of a small generating family of ultraprojective objects. For a module category, one ultraprojective generator suffices; ultraprojectivity is decomposed into projectivity and finite presentation. The proposed general structure theorem describes the resulting “algebraic” abelian topoi as categories of additive functors from $P^o$ to $(Ab)$, for a small additive category $P$. One can take $P$ essentially small and canonically equal to the full subcategory of ultraprojective objects. For $A^\wedge_{ab}$, this is the abelian Karoubi envelope of $A$, or alternatively one may use $\mathrm{Add}(A)$ for more explicit computation; the text refers back to \S\S93 and 99 (p.\ 402). The notation $P^{\&}$ and the inclusion $\varepsilon_P:P\to P^{\&}$ lead to the expected equivalence (**) between cocontinuous additive functors out of $P^{\&}$ and additive functors out of $P$. The suggested proof imitates the nonadditive representable-density argument. Its dual gives the inverse-limit version and then (***), comparing sheaves and cosheaves on the dual algebraic abelian topoi. These statements are presented as the expected additive theory, with the proof still to be checked (p.\ 403).

The next page calls an abelian topos \emph{reflexive} when it belongs to such a dual pair. Algebraicity appears sufficient, and the author wonders whether it is necessary, at least for the stronger equivalence with arbitrary additive target categories and suitable compatibility as the target varies. Remark~3 records the central limitation: for a general dual pair he does not see how to construct the comparison functor from the pairing alone. The same perplexity occurs in the nonabelian case, so the algebraic examples should not be read as an established general duality theorem (p.\ 404). Remark~4 returns to the auxiliary category $\mathcal E(A)$ from p.\ 398. Inside the functors from $B^\wedge$ to sets, it is represented by the union of the cocontinuous functors and the representable functors. Their intersection is the Karoubi envelope of $A$, equivalently the opposite of the ultraprojective objects of $B^\wedge$. The margin note describes the union by commutation with either type of small limits. These constructions are then said to extend to a small additive $P$ (pp.\ 404--405).

Case~3 resumes the intended application of (39). With an additive target $\mathcal M$, initially under the extra pseudotopos assumptions on $\mathcal M$ and its opposite, (12) and (13) yield (54), relating $\mathcal M$-valued cosheaves on $A^\wedge_{ab}$ to sheaves on $B^\wedge_{ab}$. The earlier discussion suggests that those extra assumptions should be unnecessary. The required assertion (55) says that precomposition with abelianization gives the expected equivalence for any additive $\mathcal M$ with small direct limits. Formula (56) reduces it to additive functors on $\mathrm{Add}(A)$, and (57) is the inverse-limit dual. The four-arrow diagram (58) shows that the two left comparisons are equivalent conditions because the two right comparisons are already equivalences. Finally, identifying $A^\wedge_{ab}$ with $P^{\&}$ for $P=\mathrm{Add}(A)$ reduces the remaining task to (**). The author prefers to prove that general statement and ends by saying that the suggested proof should be checked during later scratchwork. The scan's handwritten reference in the insertion is (56), and its final wording is “proof given” followed by the correction from “there” to “here” (pp.\ 405--406).


\paragraph{\S106. Dual operations, a formulaire and explicit computation (pp.\ 407--419).}
The new subsection E), dated 14 and 15 August, organizes the preceding duality into a calculus for the operations $*$ and Hom. Formula (59) relates $\mathcal M$-valued cosheaves on $A^\wedge_{ab}$ to sheaves on the opposite presheaf topos. Formula (60) instead reverses the values by passing from $\mathcal M$ to $\mathcal M^o$ while retaining the same base. Combining them gives (61), a duality expressed entirely in terms of sheaves; the author prefers this presentation for the intended homology and cohomology applications. The enlargement of diagram (58), and its dual (62), compare functors on $\mathrm{Add}(A)$, its completion under small sums, the abelianized topoi, and their opposites. The text observes that the ten displayed occurrences represent nine distinct categories, using the canonical identification $\mathrm{Add}(B)^o\simeq\mathrm{Add}(B^o)=\mathrm{Add}(A)$. The dual diagram and its introducing sentence consistently use the handwritten coefficient $\mathcal N$, an important distinction from the earlier transcription's $\mathcal M$. Formulas (63)--(64) return to the abelian pairing and its interpretation as a cosheaf (pp.\ 407--408).

Tensor structures on the two abelian topoi provide additional operations. Formula (65) is the internal Hom action of an abelian sheaf on an $\mathcal M$-valued sheaf. Through the duality, (66)--(67) express the corresponding covariant operation, denoted by a circled star. In the source this sign is repeatedly a typed star surrounded by a hand-drawn circle. The final-category example reduces these constructions to familiar Hom and tensor-product operations. Their commutation with small limits is not automatic: the text distinguishes the finitely presented projective situation from the weaker flatness detected by exactness of tensor product. The proposed converse from exactness to projectivity is explicitly rejected using flat versus free abelian groups (pp.\ 409--410). The next page retracts the initial feeling that the Hom side was already computationally understood. The data (68) are just the functors $F:A^o\to\mathcal M$, $L:A^o\to(Ab)$ and $L':B^o\to(Ab)$. Passing to the opposite coefficient category and interchanging $A,B$ produces (68') and the tautological identifications (69). Thus an adequate computation and formulaire for either operation should yield one for the other by duality; this is a program for the following pages, rather than a claim that the computation was already available (p.\ 411).

The author first broadens the Hom calculation to a topos $\mathcal A$, and provisionally to suitable pseudotopoi. The coefficient category is additive and stable under small inverse limits; the additional pseudotopos hypothesis on $\mathcal M$ is used when $\mathcal A$ is not a presheaf category. External Hom, with values in $\mathcal M$, is distinguished from internal Hom, with values in $\mathcal A_{\mathcal M}$, in (70). A sheaf $F$ is viewed as the inverse-limit-compatible functor (71). The abelianization $X\mapsto\underline Z^{(X)}$, its componentwise form for a presheaf topos, the unit $\underline Z_{\mathcal A}=\underline Z^{(e)}$ in (72), and the sections functor (73)--(74) supply the ingredients for the main list (75) (pp.\ 412--413).

Three entries of (75), circled on the page, are basic: a) evaluates external Hom on an abelianized representable; e) gives the relevant inverse-limit compatibility in the first argument; and b) is the tensor--Hom relation. The remaining entries express their internal-Hom versions, the unit identities and the relation between internal Hom and sections. The text derives them through (76), the compatibility of abelianization with products, associativity of tensor product, and specialization at the final object. The comments explicitly leave aside coherence among the canonical isomorphisms and postpone change of topos to a later section connected with homology and cohomology invariants of maps in $(Cat)$ (pp.\ 413--414).

For computation, a site $A$ presenting $\mathcal A$ allows $F$ and $L$ to be treated as the functors (77)--(79). When representables are sheaves, the canonical functor (78) is fully faithful; the text permits a change of site to reach that situation. The additive envelope and the extension $\varepsilon_{ab}:\mathrm{Add}(A)\to\mathcal A_{ab}$ in (80) turn the evaluation formula into the finite-product expression (81). The handwritten insertion and rewritten formula call the extension $\mathrm{Add}(F)^*$, although the subsequent expression (83) writes $\mathrm{Add}(F)$ without the star. A colimit presentation of $L$ by objects of $\mathrm{Add}(A)_{/L}$, compared with p.\ 403, and exactness e) then give (83). Its arrow under lim points \emph{left}: the displayed expression for $\mathrm{Hom}(L,F)$ is an inverse limit, in contrast to the arXiv witness's direct-limit arrow. The source also writes $x_i$ in the live tuple describing objects of the indexing category, then $a_i$ in the explanation; this inconsistency is preserved rather than silently repaired (pp.\ 415--416).

The remark after (83) separates the construction from the provisional assumptions. The formula makes sense for an additive $\mathcal M$ with small inverse limits, without a pseudotopos hypothesis on $\mathcal M$. This suggests an extension of the equivalence from (13), p.\ 388, when $\mathcal A$ is an actual topos. The proposed quasi-inverse still has to be checked on the abelianized representables. That check is supplied in the presheaf case by full faithfulness of (80), which gives a final object in the relevant indexing category. The author then regards the immediate computational program as complete and turns to its dual, while the proposed general extension remains qualified by the stated verification (p.\ 416).

Dualizing gives the list (84), now under the assumption $\mathcal A=A^\wedge$ and with small direct limits in $\mathcal M$. Its external star and internal circled-star formulas parallel (75), but the source retains a plain star in e'). The asymmetry of the a') formulas comes from the possible absence of binary products or sums in $A$. Formula (85) spells out the Hom calculation using the category over $a\times b$; the dual symbol $a\vee b$ initially denotes an object of $(B^\wedge)^o$ and agrees with the sum in $A$ when that sum exists. Formula (86) describes the dual indexing category concretely by triples $(x,u,v)$ with maps $a\to x$ and $b\to x$ (pp.\ 417--418). Remark~1 treats additive $A$, where finite sum and product agree. The antiequivalence (88) gives the comparison (87) between internal Hom and circled-star for $L\in\mathrm{Add}(A)$, compatible with tensor products. A faint margin note proposes extension to $\mathrm{Kar}\,\mathrm{Add}$; it is not a proof of that extension. Remark~2 begins the comparison with two more elementary operations by introducing componentwise tensor product (89); its dual is taken up on the next page (p.\ 418).

The closing page gives the dual componentwise operation (90), taking external Hom at each object and requiring small inverse limits in the additive coefficient category. Both elementary operations extend to arbitrary topoi, with abelian cosheaves replacing the opposite presheaf category in the dual construction. The author does not develop them here, but points to the role of tensor product and its derived operation for ringed topoi and anticipates a parallel role for (90). The final paragraph explains the reuse of the Hom notation: its coefficient argument lies in different categories in (65) and (90), while in the final-category case the operations agree canonically. The same reasoning applies to tensor product and the circled-star operation. This closes the notation review before the dated transition, 17 August, to ground-ring extension in \S107 (p.\ 419).
